% =====================================================================================
% REVISIÓN 1: solo los «Temas tentativos del marco teórico».
% REVISIÓN 2: el Marco teórico completo (cuando lo desarrolles, sustituye o convierte los temas tentativos en subsecciones).
% Escribe tu texto FUERA de los recuadros \begin{guia} ... \end{guia}.
% =====================================================================================

\section{Marco teórico} \label{sec:marco_teorico} % audit:req=marco_teorico

\begin{guia}[Guía: Marco teórico (Revisión 2)]
\textbf{¿Para qué sirve?} Reúne los conceptos, teorías, tecnologías, métodos y normas que se necesitan para entender y
justificar tu propuesta. No es un glosario ni un resumen de Internet: cada tema se explica \emph{en función de tu proyecto}.

\textbf{Debe responder (por cada tema):}
\begin{itemize}[nosep,leftmargin=*]
  \item ¿Qué es y cómo funciona? (definición con fuente citada y, si ayuda, un esquema o ecuación).
  \item ¿Para qué se usa en mi proyecto y por qué lo elegí frente a las alternativas?
  \item ¿Qué ventajas, limitaciones o buenas prácticas debo conocer?
\end{itemize}
Además: ¿qué metodología de desarrollo o investigación usarás y por qué? ¿qué normas o estándares aplican (calidad,
seguridad, protección de datos)?

\textbf{Organización:} una \verb|\subsection| por tema, de lo general a lo particular.

\textbf{Extensión mínima:} 400 palabras de desarrollo (sin contar los temas tentativos de la Revisión 1).

\textbf{Fuentes:} libros, artículos y documentación oficial. Cada afirmación importante lleva \verb|\cite{clave}|.

\textbf{Errores frecuentes:} temas sin relación con el proyecto; copiar y pegar; fuentes que son solo URLs de blogs; no
citar; temas que luego no se usan en el desarrollo.

\textbf{Cómo se revisa:} pertinencia, profundidad suficiente, calidad de las fuentes y que cada tema se vincule con la
propuesta.
\end{guia}

% >>> REVISIÓN 2: ESCRIBE AQUÍ el marco teórico (introducción breve y una \subsection por tema) <<<


\subsection{Temas tentativos del marco teórico} \label{sec:temas_tentativos} % audit:req=temas_marco_teorico

\begin{guia}[Guía: Temas tentativos del marco teórico (Revisión 1)]
\textbf{¿Para qué sirve?} En la primera revisión no se pide el marco teórico completo, sino la lista de los temas que
vas a desarrollar. Así el director puede confirmar que vas bien encaminado.

\textbf{Debe responder:} ¿cuáles son los temas (de 5 a 10) que necesito para sustentar el proyecto? Para cada uno, en una
o dos oraciones: ¿qué es y por qué lo necesito?

\textbf{Formato:} una lista (\verb|\begin{enumerate}| o \verb|itemize|), un tema por elemento.

\textbf{Extensión mínima:} 25 palabras. En la Revisión 2 conviertes cada tema en una subsección desarrollada.

\textbf{Cómo se revisa:} que los temas cubran lo necesario para el proyecto y no incluyan relleno.
\end{guia}

% >>> REVISIÓN 1: ESCRIBE AQUÍ la lista de temas tentativos <<<

El criterio que ordena este marco teórico es el de necesidad: se incluye únicamente aquello que el lector requiere para comprender las decisiones tomadas en el proyecto, y se excluye todo concepto que, siendo correcto, no se emplee después en el documento. Los temas se agrupan en cinco bloques que corresponden a las dimensiones del problema.

\begin{enumerate}
  \item \textbf{Gestión documental y documentos de archivo.}
  \begin{enumerate}
    \item Documento de archivo y ciclo de vida documental: producción, trámite, concentración y destino final. Sustenta por qué el sistema modela estados y no únicamente archivos almacenados.
    \item Atributos de un documento fidedigno conforme a la norma ISO 15489-1 \cite{iso15489}: autenticidad, fiabilidad, integridad y disponibilidad, y el mecanismo del sistema que garantiza cada uno.
    \item Metadatos documentales y su función en la recuperación y en la prueba.
    \item Trazabilidad y pista de auditoría: requisitos de inmutabilidad, completitud y no repudio. Es el concepto central del proyecto.
  \end{enumerate}

  \item \textbf{Ingeniería de software.}
  \begin{enumerate}
    \item Ingeniería de requerimientos conforme a la norma ISO/IEC/IEEE 29148 \cite{iso29148}: requisitos funcionales y no funcionales, elicitación, validación y trazabilidad.
    \item Modelo de calidad del producto de software de la norma ISO/IEC 25010 \cite{iso25010}, que organiza el catálogo de requisitos no funcionales.
    \item Modelos de proceso secuenciales frente a iterativos e incrementales \cite{sommerville2011, pressman2010}, y adaptación del marco de trabajo Scrum \cite{scrumguide2020} a un proyecto individual con fecha de término fija.
    \item Arquitectura de software: atributos de calidad como fuerzas de diseño, registro de decisiones arquitectónicas y arquitectura en capas \cite{bass2021}.
  \end{enumerate}

  \item \textbf{Bases de datos y modelado.}
  \begin{enumerate}
    \item Modelo relacional, normalización y modelo entidad-relación \cite{elmasri2016}.
    \item Transacciones y propiedades ACID. Sustentan la generación del folio institucional único mediante bloqueo de fila.
    \item Índices y estrategias de almacenamiento de archivos: binario en base de datos frente a sistema de archivos con referencia.
  \end{enumerate}

  \item \textbf{Tecnologías de desarrollo e infraestructura.}
  \begin{enumerate}
    \item Aplicaciones web y marco de trabajo Laravel \cite{laraveldocs}: enrutamiento, capa de acceso a datos, migraciones, validación y políticas de autorización.
    \item Sistema gestor de base de datos relacional y motor de plantillas del lado del servidor, incluyendo el control de privilegios a nivel de tabla.
    \item Virtualización y contenedores: imagen, contenedor y volumen. Sustenta la estrategia de despliegue reproducible.
    \item Administración de servidores y políticas de respaldo y recuperación.
  \end{enumerate}

  \item \textbf{Seguridad de la información.}
  \begin{enumerate}
    \item Confidencialidad, integridad y disponibilidad. Autenticación frente a autorización y control de acceso basado en roles con ámbito organizacional.
    \item Almacenamiento seguro de contraseñas mediante funciones de derivación de clave con sal.
    \item Funciones hash criptográficas aplicadas a la verificación de integridad y a la detección de cargas duplicadas.
    \item Vulnerabilidades frecuentes en aplicaciones web según el proyecto OWASP \cite{owasp2021} y cifrado en tránsito mediante TLS.
  \end{enumerate}

  \item \textbf{Marco normativo aplicable.} Disposiciones de la Ley General de Archivos \cite{lga2018} y de la Ley General de Protección de Datos Personales en Posesión de Sujetos Obligados \cite{lgpdppso2017} que inciden en el control de accesos y en el tratamiento de los datos empleados durante el desarrollo.
\end{enumerate}

Cada bloque fundamenta un elemento concreto del proyecto, de modo que ningún tema se incluye sin uso posterior: el primero sustenta los estados del documento, el versionado y la bitácora; el segundo, el catálogo de requisitos y las decisiones de arquitectura; el tercero, el modelo de datos y la generación transaccional del folio; el cuarto, la construcción de los módulos y el entorno virtualizado; el quinto, el control de acceso y las pruebas de seguridad; y el sexto, la justificación normativa de la bitácora y del tratamiento de datos de prueba.

